\section{Labels at each client}\label{sup:clients}

Table~\ref{tab:clients} gives, for each client of the two main setups,
its training questions, the questions that yield pairs and the pairs of
each preference round, and two shares among the samples that the client
draws for its own training questions in the labeling step: the samples
labeled correct, and the execution errors, the samples that fail to
execute or reach the time limit. The samples of the first preference
round come from the FedAvg SFT model after round~\setting{rounds.warmup},
those of the second from the global model after one preference round,
and on BIRD train a further set from the global model after one round
without the preference term (Section~\ref{sec:ablation} of the main
paper). At every client, one preference round raises the share of
correct samples and lowers that of execution errors, and on BIRD train
it does both by more than a round without the preference term.

\begin{table}[!ht]
\caption{Each client on its own training questions: the questions that
yield pairs and the pairs of each preference round, and the shares (\%)
of correct samples and of execution errors among the samples of the
labeling step, drawn from the FedAvg SFT model after
round~\setting{rounds.warmup} (Start), from the global model after one
preference round (After) and, on BIRD train, from the global model after
one round without the preference term (No pref.).}\label{tab:clients}
\footnotesize
%% Written by tools/paper_numbers.py. Do not edit.
\setlength{\tabcolsep}{3.5pt}%
\begin{tabular}{@{}lrrrrrrrrrrr@{}}
\toprule
 &  & \multicolumn{2}{c}{Preference round 1} & \multicolumn{2}{c}{Preference round 2} & \multicolumn{3}{c}{Correct (\%)} & \multicolumn{3}{c}{Execution errors (\%)} \\
\cmidrule(lr){3-4}\cmidrule(lr){5-6}\cmidrule(lr){7-9}\cmidrule(lr){10-12}
Client & Questions & Yield pairs & Pairs & Yield pairs & Pairs & Start & After & No pref. & Start & After & No pref. \\
\midrule
\multicolumn{12}{@{}l}{BIRD train} \\
1 & 2{,}011 & 1{,}363 & 2{,}726 & 760 & 1{,}520 & 37.9 & 59.0 & 47.7 & 31.1 & 6.3 & 20.0 \\
2 & 1{,}595 & 961 & 1{,}922 & 611 & 1{,}222 & 29.2 & 47.9 & 38.0 & 37.4 & 10.0 & 27.1 \\
3 & 2{,}080 & 1{,}379 & 2{,}758 & 874 & 1{,}748 & 33.5 & 56.4 & 43.4 & 36.8 & 8.2 & 25.7 \\
4 & 683 & 391 & 782 & 243 & 486 & 30.0 & 48.0 & 38.8 & 38.4 & 9.0 & 26.2 \\
5 & 3{,}059 & 1{,}783 & 3{,}566 & 1{,}045 & 2{,}090 & 31.3 & 51.8 & 40.1 & 29.8 & 6.3 & 20.1 \\
\midrule
\multicolumn{12}{@{}l}{Spider train} \\
1 & 895 & 457 & 914 & 167 & 334 & 74.5 & 85.0 &  & 12.7 & 4.3 &  \\
2 & 1{,}542 & 821 & 1{,}642 & 327 & 654 & 75.2 & 85.9 &  & 11.0 & 3.5 &  \\
3 & 2{,}009 & 1{,}124 & 2{,}248 & 467 & 934 & 67.9 & 81.7 &  & 14.1 & 4.3 &  \\
4 & 3{,}479 & 1{,}178 & 2{,}356 & 548 & 1{,}096 & 71.0 & 84.4 &  & 18.9 & 6.6 &  \\
5 & 734 & 433 & 866 & 168 & 336 & 71.5 & 84.4 &  & 13.7 & 5.6 &  \\
\bottomrule
\end{tabular}

\end{table}
